\uuid{LLyO}
\exo7id{1664}
\titre{exo7 1664}
\auteur{ortiz}
\organisation{exo7}
\datecreate{1999-04-01}
\isIndication{false}
\isCorrection{false}
\chapitre{Réduction d'endomorphisme, polynôme annulateur}
\sousChapitre{Diagonalisation}
\cpge{mpsi-mp2i.calcul-matriciel-systemes.e, pcsi.determinants.c, ptsi.determinants.a, mp-mpi.reduction-endomorphismes.d, pc.reduction.c, psi.reduction-endomorphismes.c, pt.reduction-endomorphismes.c, mp-mpi.reduction-endomorphismes.c, pc.reduction.b, psi.reduction-endomorphismes.b, pt.reduction-endomorphismes.b, bcpst2.valeurs-propres-vecteurs-propres.b, bcpst2.valeurs-propres-vecteurs-propres.a}
\module{Algèbre}
\niveau{L2}
\difficulte{}

\contenu{
\texte{
\'Etudier le caract\`ere diagonalisable des matrices
suivantes
}
\begin{enumerate}
    \item \question{$A=\left( \begin{array}{ccc} 3&2&4\\
-1&3&-1\\
-2&-1&-3
\end{array}
\right) \in\mathcal{M}_3(\Rr)$,}
    \item \question{$B=\left( \begin{array}{cccc}
0&\ldots &0&1\\
\vdots&&\vdots&\vdots\\
0&\ldots &0&1\\
1&\ldots&1&0
\end{array}
\right)\in\mathcal{M}_n(\Rr),n\ge 2$,}
\end{enumerate}
}
